\section{VPC - Basics}

\subsection{VPC}
\begin{itemize}
    \item A VPC must have a defined list of CIDR blocks, that cannot be changed
    \item Each CIDR within VPC: min size is /28, max size is /16 (65536 IP addresses)
    \item VPC is private, so only Private IP CIDR ranges are allowed
\end{itemize}

\subsection{Subnets}
\begin{itemize}
    \item Within a VPC, defined as a CIDR that is a subset of the VPC CIDR
    \item AAll instances within subnets get a private IP
    \item First 4 IP and last one in every subnet is reserved by AWS
\end{itemize}

\subsection{Route Tables}
\begin{itemize}
    \item Used to control where the network traffic is directed to
    \item Can be associated with specific subnets
    \item The "most specific" routing rule is always followed (192.168.0.1/24 beats 0.0.0.0/0)
\end{itemize}

\subsection{Internet Gateways IGW}
\begin{itemize}
    \item Helps a VPC connect to the internet, HA scales horizontally
    \item Acts as a NAT for instances that have a public IPv4 or public IPv6
\end{itemize}

\subsection{Public Subnets}
\begin{itemize}
    \item Has a route table that send 0.0.0.0/0 to an IGW
    \item Instances must have a public IPv4 to talk the internet
\end{itemize}

\subsection{Private Subnets}
\begin{itemize}
    \item Access internet with a NAT instance or NAT Gateway setup in a public subnet
    \item Must edit routes so that 0.0.0.0/0 routes traffic to the NAT
\end{itemize}

\subsection{VPC Basics - NAT Instance}
%TODO add diagram
\begin{itemize}
    \item EC2 instance you deploy in a public subnet
    \item Edit the route in your private subnet to route 0.0.0.0/0 to your NAT instance
    \item Not resilient to failure, limited bandwidth based on instance type, cheap
    \item Must manage failover yourself
    \item \textbf{Must disable Source/Destination Check (EC2 setting)}
\end{itemize}

\subsection{VPC Basics - NAT Gateway}
%TODO add diagram
\begin{itemize}
    \item Managed NAT solution, bandwidth scales automatically
    \item Resilient to failure within a single AZ
    \item Must deploy multiple NAT Gateways in multiple AZ for HA
    \item Has an Elastic IP, external services see the IP of the NAT gateway as the source
\end{itemize}

\subsection{Network ACL (NACL)}
\begin{itemize}
    \item Stateless firewall defined at the subnet level, applies to all instances within
    \item Support for allow and deny rules
    \item Stateless = return traffic must be explicitly allowed by rules
    \item Helpful to quickly and cheaply block specific IP addresses
\end{itemize}

\subsection{Security groups}
\begin{itemize}
    \item Applied at the instance level, only support for allow rules, no deny rules
    \item Stateful = return traffic is automatically allowed, regardless of rules
    \item Can reference other security groups in the same region (peered VPC, cross-account)
\end{itemize}

\subsection{VPC Flow Logs}
\begin{itemize}
    \item Log internet traffic going through your VPC
    \item Can be defined at the VPC level, Subnet level or ENI-level
    \item Helpful to capture "denied internet traffic"
    \item Can be sent to CloudWatch Logs and Amazon S3
\end{itemize}

\subsection{Bastion Hosts}
\begin{itemize}
    \item SSH into private EC2 instances through a public EC2 instance (bastion host)
    \item You must manage these instances yourself (failover, recovery)
    \item SSM Session Manager is a more secure way to remote control without SSH
\end{itemize}

\subsection{IPv6 in short}
\begin{itemize}
    \item All IPv6 addresses are public, total 3.4 x 10^38 addresses (vs 4.3 billion IPv4)
    \item Example CIDR:2600:LfL8:80c:a900::/56
    \item Addresses are "random" and can't be scanned online (because too many)
\end{itemize}

\subsection{VPC support for IPv6}
\begin{itemize}
    \item Create an IPv6 CIDR for VPC and use and IGW (support IPv6)
    \item Public subnet:
    \begin{itemize}
        \item Create an instance with IPv6 support
        \item Create a route table entry to ::/0 (IPv6 "all") to the IGW
        \item Public subnet (instances cannot be reached by IPv6 but can reach IPv6)
        \item Create an Egree-Only Internet Gateway in the VPC
        \item Add a route table entry for the private subnet from ::/0 to the Egress-Only IGW
    \end{itemize}
\end{itemize}


\section{VPC - Peering}
%TODO Add digram
\begin{itemize}
    \item Connect two VPC, privately using the AWS' network
    \item Make them behave as if they were in the same network
    \item Must not have overlapping CIDR
    \item VPC Peering connection is not transitive (must be established for each VPC that need to communicate with on another)
    \item You can do VPC peering with another AWS account
    \item \textbf{You must update route tables in each VPCs subnet to ensure instances can communicate}
\end{itemize}

\subsection{VPC Peering - Good to know}
\begin{itemize}
    \item VPC peering can work inter-region, cross account
    \item You can reference a security group of a peered VPC (works cross account)
\end{itemize}

\subsection{VPC Peering - Longest Prefix Match}
%TODO Add diagram
\begin{itemize}
    \item VPC uses the longest prefix match to select the most specific route
    \item Here the longest prefix for 10.0.0.77 is 10.0.0.77/32 (route table VPC A)
    - (other way of saying it is "most specific route")
\end{itemize}

\subsection{VPC Peering - Invalid Configurations}
%TODO add diagrams x 3

\subsubsection{Overlapping CIDR for IPv4}

\subsubsection{No Transitive VPC Peering}

\subsubsection{No Edge to Edge Routing}


\section{Transit Gateway}
%TODO Add digram
\begin{itemize}
    \item For having transitive peering between thousands of VPCs and on-premises, hub-and-spoke (star) connection
    \item Regional resource, can work cross-region
    \item Share cross-account using Resource Access Manager (RAM)
    \item You can peer Transit Gateways across regions
    \item Route Tables: limit which VPC can talk with other VPCs
    \item Works with Direct Connect Gateway, VPC connections
    \item Support IP Multicast (not supported by any other AWS service)
    \item Instances in a VPC can access a NAT gateway, NLB, PrivateLink and EFS in others VPCs attached to the AWS Transit Gateway
\end{itemize}

\subsection{Transit Gateway - Central NAT Gateway}
%TODO add diagram
\begin{itemize}
    \item The NAT gateway is shared in the Egress-VPC
    \item The private App VPC can access internet through the TGW
    \item In this example: the App VPCs cannot communicate with each other based on the TGW route table
\end{itemize}

\subsection{Transit Gateway - Sharing through RAM}
\begin{itemize}
    \item You can use AWS resource access manager (RAM) to share a transit gateway for VPC attachments across account or across AWS Organisation
\end{itemize}
%TODO add diagram

\subsection{Transit Gatway - Use Different Route Tables to Prevent VPC from Communicating}
%TODO add digram

\subsection{Transit Gateway to Direct Connect Gateway}
%TODO add diagram

\subsection{Transit Gateway - Inter and Intra Region Peering}
%TODO add diagram
- Billed hourly for each peering attachment, no data processing charges
- If data goes cross-region through the attachment, standard charges
- Specify the Transit Gateway ID


\section{VPC Endpoints}
%TODO add diagram
\begin{itemize}
    \item Endpoints allow you to connect to AWS Services using a private network instead of the public "www" network
    \item The scale horizontally and are redundant
    \item No more IGW, NAT, etc... to access AWS Services
    \item VPC Endpoint Gateway (S3 and DynamoDB)
    \item VPC Endpoint Interface (all services)
    \item In case of issues:
    \item Check DNS Setting Resolution in your VPC
    \item Check Route Tables
\end{itemize}

\subsection{VPC Endpoint Gateway}
\begin{itemize}
    \item Only works for S3 and DynamoDB, must create one gateway per VPC
    \item Must update route tables entries
    \item Gateway is defined at the VPC level
%TODO digagram
    \item DNS resolution must be enabled in the VPC
    \item The same public hostname for S3 can be used
    \item Gateway endpoint cannot be extended out of a VPC (VPN. DX, TGW, peering)
\end{itemize}

\subsection{VPC Endpoints Interface}
\begin{itemize}
    \item Provision an ENI that will have a private endpoint interface hostname
    \item Leverage security groups security
    \item Private DNS (setting when you create the endpoint)
    \item The public hostname of a service will resolve the private Endpoint interface hostname
    \item VPC Setting: "Enabled DNS hostnames" and "Enable DNS Support" must be "true"
    \item Example for Athena:
%        TODO add examples
    \item Interface can be accessed from Direct Connect and Site-to-Site VPN
\end{itemize}

\subsection{VPC Endpoint Policies}
\begin{itemize}
    \item Endpoint policies are JSON documents to control access to services
    \item Does not override or replace IAM user policies or service-specific policies (such as S3 bucket policies)

%TODO Json example policy

    \item Note: the IAM user can still use other SQS API from outside the VPC endpoint
    \item You could add an SQS queue policy to deny any action not done through the VPC endpoint
\end{itemize}

\subsection{VPC Endpoint Policy and S3 bucket policy}
\begin{itemize}
    \item VPC Endpoint Policy to restrict access to bucket "my_secure_bucket"
%TODO add JSON policy

    \item VPC Endpoint to allow access to Amazon Linux 2 repositories
%TODO add JSON policy

    \item S3 bucket policy may have
    \begin{itemize}
        \item Condition: "aws:sourceVpce":"vpce-la2b3c4d" to Deny any traffic that doesn't come from a specific VPC endpoint (more secure)
        \item Condition: "aws:sourceVpc":"vpc-111bbb22" for a specific VPC
        \item The aws:sourceVpc condition only works for VPC endpoints in case you have multiple endpoints and want to manage access to your S3 buckets for all yyou endpoints
        \item The S3 bucket policies policies can restrict access only from a specific public IP address or an elastic IP address. \textbf{You can't restrict based on private IP}
        \item Therefore aws:Source condition doesn't apply for VPC endpoints
    \end{itemize}
\end{itemize}

\subsubsection{Example S3 bucket policies}
\begin{itemize}
    \item S3 bucket policy to restrict to specific VPC endpoint
%TODO example JSON policy
    \item S3 bucket policy to restrict to an entire VPC (multiple VPC Endpoints)
%TODO example JSON policy
\end{itemize}

\subsection{VPC Endpoint Policies for S3 Troubleshooting}
%TODO diagram


\section{PrivateLink}
\begin{itemize}
    \item Most secure and scalable way to expose a service to 1000s of VPC (own or other accounts)
    \item Does not require VPC peering, internet gateway NAT, route tables...
    \item Requires a network load balancer (Service VPC) and ENI (CustomerVPC)
    \item If the Network Load Balancer (NLB) is in multiple AZ and the ENI in multiple AZ the solution is fault tolerant (good)
\end{itemize}

%TODO diagram

\subsection{PrivateLink for Amazon S3 with Direct Connect}
%TODO diagram

\subsection{VPC Endpoints / PrivateLink and VPC peering}
%TODO diagram

\subsection{Site to Site VPN (AWS Managed VPN)}
\begin{itemize}
    \item on-premises
    \begin{itemize}
        \item setup a software of hardware VPN appliance to your on-premises network
        \item The on-premises VPN should be accessible using a public IP
    \end{itemize}
    \item AWS-side
    \begin{itemize}
        \item Setup a Virtual Private Gateway (VGW) and attach to you VPC
        \item Setup a Customer Gateway to point the on-premises VPN appliance
        \item Two VPN connections (tunnels) are created for redundancy,encrypted using IPSec
        \item Can optionally accelerate it using Global Accelerator (for worldwide networks)
    \end{itemize}
\end{itemize}
%TODO diagram

\subsection{Route Propagation in Site-to-Site VPN}
%TODO diagram

\subsubsection{Static Routing}
\item Create static route in corporate data centre for 10.0.0.1/24 through the CGW
\item Create static route in AWS for 10.3.0.0/20 through the VGW
\item Dynamic Routing (BGP)
\item Using BGP (Border Gateway Protocol) to share routes automatically (eBGP for internet)
\item We don't need to update the routing tables, it will be done for us dynamically
\item Just need to specify the ASN (Autonomous System Number) of the CGW and VGW

\subsection{Site to Site VPN and Internet Access}
\begin{itemize}
    \item Not Okay (blocked by NAT Gateway restrictions)
%TODO add diagram
    \item Okay (self manage NAT instance - more control)
%TODO add diagram
    \item Okay (alternative to NAT instances / gateway)
%TODO add diagram
\end{itemize}

\subsection{AWS VPN CloudHub}
%TODO add diagram
\begin{itemize}
    \item Can connect up to 10 customer gateway for each Virtual Private Gateway (VGW)
    \item Low cost hub-and-spoke model for primary or secondary network connectivity between locations
    \item Provide secure communication between sites, if you have multiple VPN connections
    \item It's a VPN connection so it goes over the public internet
    \item Can be a failover connection between your on-premises locations
\end{itemize}

\subsection{VPN to multiple VPC}
%TODO add diagram
\begin{itemize}
    \item For VPN-based customers, AWS recommends creating a separate VPN connection for each customer VPC
    \item Direct Connect is recommended because it has Direct Connect Gateway
\end{itemize}

\subsection{Shared Services VPN}
%TODO add diagram
\begin{itemize}
    \item Create a VPN connection between on-premises and shared service VPC
    \item Replicate services, applications, databases between on-premises and the Shared Services VPC or deploy proxies in the shared service VPC
    \item Do VPC peering between the VPC and the shared service VPC
    \item VPCs can directly access the Shared Service VPC services and do not need VPN connections to on-premises
\end{itemize}


\section{AWS S2S VPN}


\section{AWS Client VPN}
%TODO add diagram
\begin{itemize}
    \item Connect from your computer using OpenVPN to your private network in AWS and on-premises
\end{itemize}

\subsection{Client VPN - Peered VPC}
\begin{itemize}
    \item Client VPN is compatible with VPC peering
\end{itemize}
%TODO add diagram

\subsection{Client VPN - Access to On-Premises}
\begin{itemize}
    \item Access On-Premises resources through AWS with ClientVPN
\end{itemize}
%TODO add diagram

\subsection{Client VPN - Internet Access}
%TODO add diagram

\subsection{Client VPN - Transit Gateway}
%TODO add diagram x 2


\section{Direct Connect}
\begin{itemize}
    \item Provides a dedicated private connection from a remote network to your VPC
    \item Dedicated connection must be setup between you DC and AWS Direct Connection location
    \item More expensive than running a VPN solution
    \item Private access to AWS service through VIF
    \item Bypass ISP, reduce network cost, increase bandwidth and stability
    \item Not redundant by default (must setup a failover DX or VPN)
\end{itemize}

\subsection{Direct Connect - Virtual Interfaces (VIF)}
\begin{itemize}
    \item Public VIF - connect to Public AWS Endpoints (S3 buckets, EC2 service, anything AWS ...)
    \item Private VIF - connect to resources in your VPC (EC2 instances, ALB, ...)
    \item Transit Virtual Interface - connect to resource in a VPC using a Transit Gateway
    \item \underline{VPC Interface Endpoints can be accessed through Private VIF}
\end{itemize}

\subsection{Direct Connect Diagram}
%TODO add direct connect diagram

\subsection{Direct Connect - Connection Types}
\begin{itemize}
    \item Dedicated Connections: 1 Gbps, 10 Gbps, 100 Gbps capacity
    \item Physical ethernet port dedicated to a customer
    \item Request made to AWS first, the completed by AWS Direct Connect Partners
    \item Hosted Connections: 50Mbps, 500Mbps to 10Gbps
    \item Connection requests are made via AWS Direct Connect Partners
    \item Capacity can be added or removed on demand
    \item 1,2,5,10 Gbps available at select AWS Direct Connect Partners
    \item Lead times are often longer than 1 month to establish a new connection
\end{itemize}

\subsection{Direct Connect - Encryption}
%TODO add diagram
\begin{itemize}
    \item Data in transit is not encrypted but is private
    \item AWS Direct Connect + VPN provides an IPsec-encrypted private connection
    \item VPN over Direct Connect connection Uses Public VIF
    \item Good for an extra level of security, but slightly more complex to put in place
\end{itemize}

\subsection{Direct Connect - Link Aggregation Groups (LAG)}
%TODO add diagram
\begin{itemize}
    \item Get increased speed and failover by summing existing DX connections in a logical one
    \item Can aggregate up to 4 connections (active-active mode)
    \item Can add connections over time to the LAG
    \item All connection in the LAG:
    \item Must be dedicated connections
    \item Must have the same bandwidth
    \item Must terminate at the same AWS Direct Connect Endpoint
    \item Can set a minimum number of connections for the LAG to function
\end{itemize}

\subsection{Direct Connect Gateway}
\begin{itemize}
    \item If you want to setup a Direct Connect to one or more VPC in many different regions (same/ cross account), you must use a Direct Connect Gateway
\end{itemize}
%TODO add diagram

\subsection{Site-to-site Active-Active Connection}
%TODO Active-Active VPN connection diagram

\subsection{Direct Connect - High Availability}
%TODO Multiple connections at multiple AWS direct connect location diagram
%TODO Backup VPN connection diagram

\subsection{Direct Connect Gateway - SiteLink}
%TODO add diagram
\begin{itemize}
    \item Allows you to send data from on Direct Connect location to another
    \item Bypassing AWS regions
    \item Use cases: create private network connections between on-premises data centres by connecting them to Direct Connect locations
    \item Data is sent over the fastest path between Direct Connect locations
\end{itemize}


\section{On-Premise Redundant Connections}
%TODO ??????


\section{VPC Flow Logs}
\begin{itemize}
    \item Capture information about IP traffic going into your interfaces:
    \item VPC flow logs
    \item Subnet flow logs
    \item Elastic Network Interface (ENI) Flow Logs
    \item Help to monitor and troubleshoot connectivity issues
    \item Flow logs data can go to S3, CloudWatch Logs and Kinesis Data Firehose
    \item Captures network information from AWS managed interfaces too: ELB, RDS, ElastiCache, Redshift, WorkSpaces, NATGW, Transit Gateway.....
\end{itemize}

\subsection{VPC Flow Logs Syntax}
\begin{itemize}
    \item srcaddr and dstaddr - help identify problematic IP
    \item srcport and dstport - help identify problematics ports
    \item Action - success or failure of the request due to Security Group / NACL
    \item Can be used for analytics on usage patterns, or malicious behaviour
    \item Query VPC flow logs using Athena on S3 or CloudWatch logs Insights
    \item Flow logs examples: https://docs.aws.amazon.com/vpc/latest/userguide/flow-logsrecords-
    examples.html
\end{itemize}

\subsection{VPC Flow Logs - Troubleshoot SG and NACL issues}
Look at the "Action" field

\subsubsection{Incoming Requests}
\begin{itemize}
    \item Inbound REJECT => NACL or SG
    \item Inbound ACCEPT, Outbound Reject => NACL
\end{itemize}
%TODO add diagram

\subsubsection{Outgoing Requests}
\begin{itemize}
    \item Outbound REJECT => NACL or SG
    \item Outbound ACCEPT, Inbound REJECT => NACL
\end{itemize}
%TODO add diagram

\subsection{VPC Flow Logs - Architectures}
%TODO add diagram of VPC Flow Logs - Architectures

\subsection{VPC Flow Logs - with NAT Gateway}
%TODO add diagram
\begin{itemize}
    \item My Virtual Private Cloud (VPC) flow logs show Action = ACCEPT for inbound traffic coming from public IP addresses.
    However, my understand of network address translation (NAT) gateway was that they don't accept traffic from the internet.
    Is my NAT gateway accepting inbound traffic from the internet
    \item Inbound traffic is permitted by Security Group or NACLs
    \item Traffic? isn't permitted by the NAT gateway, it's dropped
    \item To confirm run the following query in CloudWatch Log Group
    \item Make 'xxx.xxx' the first two octets of your VPC CIDR
    \item Replace Public IP with the IP you see in logs
    \item You will see traffic on the Private IP of the NAT Gateway but nowhere else: traffic was unsolicited and then dropped
\end{itemize}


\section{AWS Network Firewall}

\subsection{Network Protection on AWS}
\begin{itemize}
    \item To protect network on AWS, we've seen
    \item Network Access Control Lists (NACLs)
    \item Amazon VPC security groups
    \item AWS WAF (protect against malicious requests)
    \item AWS Shield and AWS Shield Advanced
    \item AWS Firewall Manager (to manage them across accounts)
    \item But what if we want to protect in a sophisticated way out entire VPC?
\end{itemize}

\subsection{AWS Network Firewall}
%TODO add diagram
\begin{itemize}
    \item Protected your entire Amazon VPC
    \item From Layer 3 to Layer 7 protection
    \item Any direction, you can inspect
    \item VPC to VPC traffic
    \item Outbound to internet
    \item Inbound from internet
    \item To / from Direct Connect and Site-to-Site VPN
    \item Internally, the AWS Network Firewall uses the AWS Gateway Load Balancer
    \item Rules can be centrally managed cross-account by AWS Firewall Manager to apply to many VPCs
\end{itemize}

\subsection{Network Firewall - Fine Grained Controls}
\begin{itemize}
    \item Supports 1000s of rules
    \item IP and port - example: 10,000s of IPs filtering
    \item Protocol - example: block the SMB protocol for outbound communications
    \item Stateful domain list rule groups, only allow outbound traffic to *.mycorp.com or third-party software repo
    \item General pattern matching using regex
    \item Traffic filtering: Allow drop or alert for the traffic that matches the rules
    \item Active flow inspection to protect against network threats with intrusion-prevention capabilities (like Gateway Load Balancer, but all managed by AWS)
    \item Send logs or rule matches to Amazon S3, CloudWatch Logs, Kinesis Data Firehose
\end{itemize}

\subsection{AWS Network Firewall Architecture}
%TODO Add diagram of AWS network Firewall - Architecture
